\documentclass[11pt,a4paper]{article}
\usepackage[T1]{fontenc}
\usepackage[utf8]{inputenc}
\usepackage{amsmath,amssymb,amsthm}
\usepackage[margin=1in]{geometry}
\usepackage[colorlinks=true,linkcolor=blue,urlcolor=blue]{hyperref}
\theoremstyle{plain}
\newtheorem{theorem}{Theorem}
\newtheorem{lemma}[theorem]{Lemma}
\newtheorem{proposition}[theorem]{Proposition}
\newtheorem{corollary}[theorem]{Corollary}
\theoremstyle{definition}
\newtheorem{remark}[theorem]{Remark}

\title{A lattice count, not a walk count: OEIS A225315}
\author{Adrian Perez Fontelles\\ \small Independent researcher}
\date{7 September 2026}
\begin{document}
\maketitle

\begin{abstract}
OEIS A225315 carries an empirical recurrence of order $36$ contributed by R.~H.~Hardin. It is
true, and it is decidable rather than empirical. A nondecreasing row of $-1$s and $1$s is
determined by one number --- how many $-1$s it holds --- so the entry's condition collapses to
a single linear equation on a box of lattice points, and the arrays disappear entirely. The
count is then the number of lattice points in a dilate of a fixed rational polytope, hence a
quasi-polynomial in $n$; its degree and its period can both be named, which bounds by
$S=13440$ the order of a linear recurrence it must satisfy. Whether the conjectured recurrence
is that recurrence is then a finite exact computation. There is no digraph and no matrix
anywhere in this argument.
\end{abstract}

\noindent\small 2020 Mathematics Subject Classification. 05A15, 05B45, 15B36.\normalsize

\section{The conjecture}

OEIS A225315 is ``Number of 8Xn -1,1 arrays such that the sum over i=1..8,j=1..n of i*x(i,j) is zero and rows are nondecreasing (ways to put n thrusters pointing east or west at each of 8 positions 1..n distance from the hinge of a south-pointing gate without turning the gate).''. It has offset $1$ and begins
\[
14, 217, 1576, 7445, 26498, 77721, 197440, 449693,\ \dots
\]
The entry states, as a conjecture contributed by its author and never marked settled:
\begin{quote}
Empirical: a(n) = a(n-1) +a(n-2) -a(n-5) -a(n-7) -a(n-9) +a(n-11) +2*a(n-12) +a(n-13) +a(n-15) -a(n-16) -a(n-17) -2*a(n-18) -a(n-19) -a(n-20) +a(n-21) +a(n-23) +2*a(n-24) +a(n-25) -a(n-27) -a(n-29) -a(n-31) +a(n-34) +a(n-35) -a(n-36)
\end{quote}
As of the ``Last modified'' line on the live entry (Oct 03 2025 23:43:12, revision 8) this is still
recorded as empirical, and nothing on the entry records it as proved.

\section{The array disappears}

A row of $-1$s and $1$s that is nondecreasing is a block of $-1$s followed by a block of $1$s, so it is determined by a single number: how many $-1$s it holds. Write $k_i\in\{0,1,\dots,n\}$ for that number in row $i$. Then row $i$ sums to $(n-k_i)-k_i = n-2k_i$, and the entry's condition
\[ \sum_{i=1}^{H}\sum_{j=1}^{n} i\,x(i,j) \;=\; \sum_{i=1}^{H} i\,(n-2k_i) \;=\; 0 \]
becomes $\sum_{i=1}^{8} i\,k_i = T(n)$ with $T(n)=18n$. The arrays are gone: what is left is the lattice points of the box $\{0,\dots,n\}^{8}$ on one hyperplane.

\section{Why that is enough}

Those lattice points are the ones in the $n$-th dilate of a fixed rational polytope, so the count is a quasi-polynomial in $n$. That is the whole reason a linear recurrence can be proved here rather than observed: a quasi-polynomial of degree $d$ and period $P$ satisfies the recurrence whose characteristic polynomial is $(x^{P}-1)^{d+1}$, so it is $C$-finite of order at most $P(d+1)$.

Both numbers can be named. By inclusion-exclusion over which coordinates exceed $n$,
\[ a(n)\;=\;\sum_{S\subseteq\{1,\dots,H\}}(-1)^{|S|}\,D\bigl(T(n)-(n+1)\sigma(S)\bigr), \]
where $\sigma(S)$ is the sum of $S$ and $D(m)$ counts the solutions of $\sum i\,k_i=m$ in non-negative integers with no upper bound --- the classical denumerant for the parts $1,\dots,8$. $D$ is a quasi-polynomial in $m$ of degree $7$ with period dividing $\mathrm{lcm}(1,\dots,8)=840$; each argument above is an integer-linear function of $n$, which preserves the degree and can only shrink the period, and the parity case at worst doubles it. So $a$ is a quasi-polynomial of degree at most $7$ and period dividing $1680$, hence satisfies a linear recurrence of order at most $S=13440$.

That bound is the certificate. It is also what makes the computation cheap: $D$ is one table, so thousands of exact terms of $a$ cost almost nothing.
so $a$ is $C$-finite of order at most $S=13440$, and the alignment between the model and the
entry's own $n$ is fixed by matching against every published term.

\section{The proof}

\begin{lemma}\label{lem:crit}
Let $q$ be the characteristic polynomial of the conjectured recurrence, of order $36$,
and let $u_j=a(j+36)-\sum_i c_i\,a(j+36-i)$ be its residual. Then $u$ satisfies the
same linear recurrence of order $S=13440$ that $a$ does, so if $u_j=0$ for $S$ consecutive
indices then $u_j=0$ for every larger $j$, and the last nonzero $u_j$ pins the threshold
exactly.
\end{lemma}

\begin{proof}
A sequence annihilated by a linear operator with constant coefficients has every linear
combination of its shifts annihilated by the same operator, and $u$ is such a combination of
shifts of $a$. A sequence satisfying a linear recurrence of order $S$ and vanishing at $S$
consecutive indices vanishes from the first of them onwards, since every later value is a
fixed combination of the $S$ before it.
\end{proof}

\begin{theorem}
\[
a(n)\;=\;a(n-1) + a(n-2) - a(n-5) - a(n-7) - a(n-9) + a(n-11) + 2\,a(n-12) + a(n-13) + a(n-15) - a(n-16) - a(n-17) - 2\,a(n-18) - a(n-19) - a(n-20) + a(n-21) + a(n-23) + 2\,a(n-24) + a(n-25) - a(n-27) - a(n-29) - a(n-31) + a(n-34) + a(n-35) - a(n-36)
\]
for every $n>35$.
\end{theorem}

\begin{proof}
The terms $a(n)$ were computed exactly from the lattice-point formula, far enough to cover
$S=13440$ consecutive residuals, and the residuals were formed in exact integer arithmetic.
They vanish from the index corresponding to $n=35+1$ onwards and go on vanishing for more
than $S$ consecutive indices, which by Lemma~\ref{lem:crit} settles every larger $n$ at
once.
\end{proof}

\section{Verification}

Four checks, each independent of the algebra above.

First, the model reproduces all $26$ terms the entry publishes, exactly, in integer
arithmetic. This is what ties the reduction to the entry: the lattice count is derived by reading the
entry's English, and a misreading gives different counts.

Second, the reading was pinned before the reduction was trusted, by a program that enumerates
the arrays themselves, one by one, from the entry's own words --- every nondecreasing
row of plus and minus ones, the weighted sum computed directly --- and compares the count
with the lattice-point formula. That program shares no code with the lattice-point formula, so an error in one does not
hide an error in the other.

Third, the conjectured recurrence was evaluated directly on the published terms, with
nothing from the argument above involved, and holds wherever the proved range and the data
overlap.

Fourth, the residual test of Lemma~\ref{lem:crit} was carried out in exact integer
arithmetic over every index up to the certified bound, not on a sampled prefix.

\begin{thebibliography}{9}
\bibitem{oeis} The OEIS Foundation, \emph{The On-Line Encyclopedia of Integer
Sequences}, \url{https://oeis.org}, 2026. Sequence A225315.
\bibitem{stanley} R.~P.~Stanley, \emph{Enumerative Combinatorics}, Volume 1, 2nd ed.,
Cambridge University Press, 2011. (Quasi-polynomials and denumerants, Sections 4.4 and 4.6;
lattice points in dilates of rational polytopes, Section 4.6.)
\bibitem{fs} P.~Flajolet and R.~Sedgewick, \emph{Analytic Combinatorics}, Cambridge
University Press, 2009.
\bibitem{beck} M.~Beck and S.~Robins, \emph{Computing the Continuous Discretely}, 2nd
ed., Springer, 2015. (Ehrhart quasi-polynomials; the denumerant.)
\end{thebibliography}

\end{document}
